\section{Electron Fake Factor Measurement}%
\label{sec:bkg-fakes-electron}

Fake electrons arise from three background sources. They are either
mis-reconstructed jets, mostly originating from light quark or gluon
fragmentation, photon conversions or real electrons arising from secondary
decays of light and heavy flavour mesons within jets. Fake electrons are
normally less isolated than prompt electrons and have a lower chance of passing
tighter identification levels. The simulation does not offer an optimal
modelling of fake electrons which ultimately arise as the final product of
a chain prompted by strong interactions.

The electron fake contribution to the measurement region is estimated as described
in \cref{sec:bkg-fakefactor}, which involves a data driven estimate of
the fake factor.

Electron fake factors are measured by selecting a region in data that
is significantly enriched in fake electrons while still kinematically close
enough to the measurement region. This is not easily achieved because fake electrons
can very closely mimic the attributes of prompt electrons, and the contribution
of prompt leptons cannot be completely suppressed, so one needs to explicitly
subtract the expected prompt electron contamination using MC predictions.


\subsection{Electron Fake Factor Measurement}

To select a region with a significant content of fake electrons, a special
single electron region is defined. The cut on \( \met + \mtW \) has been
reversed to require the observable smaller than \qty{60}{\GeV}. Additionally,
to increase the statistics, jet and \bjet selections have been loosened
to at least 3 jets of which at least one should be a \bjet. The region requirements
are summarised in \cref{tab:FakeRegionDefinitions}.

\begin{table}[htbp]
  \begin{center}
    \begin{tabular}{l c c}
      \toprule
      Selection            & Electron fakes region             & Muon fakes region  \\
      \midrule
      \( N \)(central jet) & \multicolumn{2}{c}{\( \geq 3\)}                        \\
      \( N_{\bjets}\)      & \multicolumn{2}{c}{\( \geq 1 \)}                       \\
      \( \met + \mtW \)    & \multicolumn{2}{c}{\( (0, 60) \)}                      \\
      \( N \)(lep)         & 1                                 & 1                  \\
      \( N(\ell) \)        & 1                                 & 0                  \\
      \( N(\mu) \)         & 0                                 & 1                  \\
      \( \pt(\ell) \)      & \( (30, \infty) \)                & \( (30, \infty) \)                \\
      \( \pt(j1)\)         & ---                               & \( (35, \infty) \) \\
      \( \Delta \phi(\mu, \mathrm{jet})\) & --- & \( \geq 2.7\) \\
      \bottomrule
    \end{tabular}
  \end{center}
  \caption{A summary of the event selection for fake measurement regions.}%
  \label{tab:FakeRegionDefinitions}
\end{table}

All object definitions are the same as defined in \cref{ch:definitions}.
The defined region is called the \textit{fakes-enriched region}.
Two sub-regions are defined out of the fakes-enriched region, \textit{tight}
and \textit{loose}, which are filled with electrons satisfying the corresponding
tightness requirement.

\begin{figure}[htbp]
  \centering
  \includegraphics[width=0.4\textwidth]{fakes/fake_factors_e_nominal}%
  \label{fig:Fakes_Electrons-nominal}
  \caption{Electron fake factors measured as a function of \pt for different \eta bins.
  }%
  \label{fig:Fakes_Electrons}
\end{figure}

\begin{figure}[htbp]
  \centering
  \subfloat[]{
    \includegraphics[page=7, width=0.4\textwidth]{fakes/Fakes_Electrons_e_nominal}%
    \label{fig:Fakes_Electrons-pt-loose}
  }
  \subfloat[]{
    \includegraphics[page=3, width=0.4\textwidth]{fakes/Fakes_Electrons_e_nominal}%
    \label{fig:Fakes_Electrons-pt-tight}
  }
  \\
  \subfloat[]{
    \includegraphics[page=12, width=0.4\textwidth]{fakes/Fakes_Electrons_e_nominal}%
    \label{fig:Fakes_Electrons-eta-loose}
  }
  \subfloat[]{
    \includegraphics[page=11, width=0.4\textwidth]{fakes/Fakes_Electrons_e_nominal}%
    \label{fig:Fakes_Electrons-eta-tight}
  }

  \caption{Fakes-enriched regions in the nominal selection.
    \protect\subref{fig:Fakes_Electrons-pt-loose} \pt distribution of loose electrons in the fakes-enriched region,
    \protect\subref{fig:Fakes_Electrons-pt-tight} \pt distribution of tight electrons in the fakes-enriched region,
    \protect\subref{fig:Fakes_Electrons-eta-loose} \eta distribution of loose electrons in the fakes-enriched region,
    \protect\subref{fig:Fakes_Electrons-eta-tight} \eta distribution of tight electrons in the fakes-enriched region.
    All the distributions show data events and the prompt MC component subtracted from data, to ensure a fake dominated region.
  }%
  \label{fig:Fakes_Electrons-distributions}
\end{figure}

The fakes-enriched regions are shown in \cref{fig:Fakes_Electrons-distributions}.
Selected events still contain prompt electrons after all selection requirements
are applied. These residual prompt electrons are subtracted from data before
the calculation of fake factors using the Monte Carlo simulation.
The MC samples used are \Wjets (before the fit), \Zjets, \ttbar and single top, diboson, and rare
top decays. The MC subtraction is much larger in the tight region
compared to the loose region and the number of subtracted prompt leptons
in that region amounts to up to \qty{100}{\%} of all electrons.
The \pt dependence of the fake factor is shown in \cref{fig:Fakes_Electrons}.
The fake factor was measured in three \eta slices (two slices for the forward
region and one slice for the barrel region). The last \pt bin includes the
overflow and is extended to infinity and no extrapolation of the fake
factor is performed.


\subsection{Systematic Uncertainties on Electron Fake-factors}%
\label{sec:bkg-fakes-electron-sys}

The main sources of the systematic uncertainties of the fake factor arise from
Monte Carlo modelling of the subtracted leptons, from different composition
of fake electrons in the fake enriched region compared to the measurement region,
and from the normalisation of Monte Carlo samples in the fakes-enriched region.

Given that the amount of tight fakes is very low, only the MC normalisation
uncertainty is evaluated by varying all samples simultaneously by \qty{10}{\%} to account
for cross-section and luminosity uncertainties as well as the modelling of
the prompt MC used in the subtraction procedure.

Additionally the effect of the \( \met + \mtW \) cut is tested by varying it
to \qty{50}{\GeV} and \qty{70}{\GeV}. The effect of this variation on the fake
factors is negligible.

The effect on the number of \(b-\)jets was tested by varying \( N_{\bjets}\) to include at least 2 \(b-\)jets. This effect isn't negligible and can be quite large in some bins in the electron \pt distribution. 

This region is kinematically quite far from the measurement region so the observed
variation in the fake factor due to varying composition of the fakes enriched
sample is a quite conservative estimate of the systematic uncertainty.

% \begin{table}[htbp]
%   \begin{center}
%     \begin{tabular}{ c c }
%       \toprule
%       Variation & Purpose \\
%       \midrule
%       Flipped requirement on number of jets      & fakes composition                 \\
%       Removed \met requirement      & fakes composition                 \\
%       MC scaled up by \qty{10}{\%}     & MC modelling, cross-section and luminosity \\
%       MC scaled down by \qty{10}{\%}   & MC modelling, cross-section and luminosity \\
%       \bottomrule
%     \end{tabular}
%   \end{center}

%   \caption{Summary of the variations used for the determination of the systematic uncertainty of the fake factors.}%
%   \label{tab:fakes-electron-sys}
% \end{table}

The resulting fake factors are shown on \crefrange{fig:Fakes_Electrons-elsys1}{fig:Fakes_Electrons-elsys4}.
Combined systematics of the fake factors was calculated by adding all variations
and the statistical uncertainty in quadrature. The measurement is relatively
stable overall up to very high electron \pt but the maximum uncertainty can get
very high, up to \qty{100}{\%}.

\begin{figure}[htbp]
  \centering
  \subfloat[]{
    \includegraphics[page=1, width=0.4\textwidth]{fakes/fake_factors_e_final}%
    \label{fig:Fakes_Electrons-elsys1}
  }
  \subfloat[]{
    \includegraphics[page=2, width=0.4\textwidth]{fakes/fake_factors_e_final}%
    \label{fig:Fakes_Electrons-elsys2}
  }
  \\
  \subfloat[]{
    \includegraphics[page=3, width=0.4\textwidth]{fakes/fake_factors_e_final}%
    \label{fig:Fakes_Electrons-elsys3}
  }
  \subfloat[]{
    \includegraphics[page=4, width=0.4\textwidth]{fakes/fake_factors_e_final}%
    \label{fig:Fakes_Electrons-elsys4}
  }

  \caption{The measured fake factor with systematic uncertainties.
    \protect\subref{fig:Fakes_Electrons-elsys1} fake factor variations for the first \eta bin,
    \protect\subref{fig:Fakes_Electrons-elsys2} fake factor variations for the second \eta bin,
    \protect\subref{fig:Fakes_Electrons-elsys3} fake factor variations for the third \eta bin and
    \protect\subref{fig:Fakes_Electrons-elsys4} fake factor variations for the fourth \eta bin.
  }
\end{figure}


\subsection{Validation of Fake-factors for Electrons}%
\label{subsec:bkg-fakes-electron-validation}

The measured fake factors were validated by performing the direct closure test
in the fakes-enriched region. The same event selection was used as for the
estimation, but no special binning has been applied. Fakes have been estimated
from data using the measured fake factors. The reconstructed object selection
was the same as the nominal event selection (\cref{ch:definitions}).

Prompt Monte Carlo processes that were used are the same as for the estimation.
The remaining events are described well by the fake factor method as seen on
\cref{fig:Fakes_ClosureElectrons}.
The only considered systematic uncertainty in this region was
the systematic uncertainty of the fake factor. Overall the agreement is good
even up to very high electron energies.

\begin{figure}[htbp]
  \centering
  \subfloat[]{
    \includegraphics[page=2, width=0.4\textwidth]{fakes/Fakes_ClosureElectrons_e_validation_sys}%
    \label{fig:Fakes_ClosureElectrons-pt}
  }
  \subfloat[]{
    \includegraphics[page=5, width=0.4\textwidth]{fakes/Fakes_ClosureElectrons_e_validation_sys}%
    \label{fig:Fakes_ClosureElectrons-eta}
  }
  \\
  \subfloat[]{
    \includegraphics[page=7, width=0.4\textwidth]{fakes/Fakes_ClosureElectrons_e_validation_sys}%
    \label{fig:Fakes_ClosureElectrons-met}
  }
  \subfloat[]{
    \includegraphics[page=13, width=0.4\textwidth]{fakes/Fakes_ClosureElectrons_e_validation_sys}%
    \label{fig:Fakes_ClosureElectrons-mt}
  }

  \caption{The electron fake-factors closure.
    \protect\subref{fig:Fakes_ClosureElectrons-pt} electron \pt distribution,
    \protect\subref{fig:Fakes_ClosureElectrons-eta} electron \eta distribution,
    \protect\subref{fig:Fakes_ClosureElectrons-met} \met distribution, and
    \protect\subref{fig:Fakes_ClosureElectrons-mt} \mtW distribution.
    The dashed band shows the systematic uncertainty associated to the fake factor measurement.
  }%
  \label{fig:Fakes_ClosureElectrons}
\end{figure}
